\documentclass{article}

% packages
\usepackage{amsmath, amsthm, thmtools, amsfonts, amssymb, luacode, catchfile, tikzducks, hyperref, ifthen}
\ifcsname c@kobocompile\endcsname
	\usepackage[a5paper, total={1072pt, 1448pt}, margin=10pt, includeheadfoot]{geometry} % set page margins
\else
	\usepackage[a4paper, margin=50pt, includeheadfoot]{geometry}
\fi
\usepackage[shortlabels]{enumitem}
\usepackage[skip=3pt, indent=0pt]{parskip}

% language
\usepackage[bidi=basic, layout=tabular, provide=*]{babel}
\ifcsname c@english\endcsname
	\babelprovide[main, import]{english}
\else
	\babelprovide[main, import]{hebrew}
	\babelprovide{rl}
\fi
%\babelfont{rm}{Libertinus Serif}
\babelfont{rm}[Renderer=Harfbuzz]{Libertinus Serif}
\babelfont{sf}{Libertinus Sans}
\babelfont{tt}{Libertinus Mono}

% style
\AddToHook{cmd/section/before}{\clearpage}	% Add line break before section
\linespread{1.3}
\setcounter{secnumdepth}{0}		% Remove default number tags from sections, this won't do well with theorems
\AtBeginDocument{\setlength{\belowdisplayskip}{3pt}}
\AtBeginDocument{\setlength{\abovedisplayskip}{3pt}}
\graphicspath{ {../images/} }

% operators
\DeclareMathOperator\cis{cis}
\DeclareMathOperator\Sp{Sp}
\DeclareMathOperator\tr{tr}
\DeclareMathOperator\im{Im}
\DeclareMathOperator\re{Re}
\DeclareMathOperator\diag{diag}
\DeclareMathOperator*\lowlim{\underline{lim}}
\DeclareMathOperator*\uplim{\overline{lim}}
\DeclareMathOperator\rng{rng}
\DeclareMathOperator\Sym{Sym}
\DeclareMathOperator\Arg{Arg}
\DeclareMathOperator\Log{Log}
\DeclareMathOperator\dom{dom}
\DeclareMathOperator\supp{Supp}
\DeclareMathOperator\var{Var}
\DeclareMathOperator\cov{Cov}

% commands
%\renewcommand\qedsymbol{\textbf{מש''ל}}
%\renewcommand\qedsymbol{\fbox{\emoji{lizard}}}
\newcommand{\Aa}[0]{\mathcal{A}}
\newcommand{\Bb}[0]{\mathcal{B}}
\newcommand{\CC}[0]{\mathbb{C}}
\newcommand{\Cc}[0]{\mathcal{C}}
\newcommand{\EE}[0]{\mathbb{E}}
\newcommand{\FF}[0]{\mathbb{F}}
\newcommand{\Ff}[0]{\mathcal{F}}
\newcommand{\Ii}[0]{\mathcal{I}}
\newcommand{\Gg}[0]{\mathcal{G}}
\newcommand{\Ll}[0]{\mathcal{L}}
\newcommand{\Mm}[0]{\mathcal{M}}
\newcommand{\NN}[0]{\mathbb{N}}
\newcommand{\Nn}[0]{\mathcal{N}}
\newcommand{\PP}[0]{\mathbb{P}}
\newcommand{\Pp}[0]{\mathcal{P}}
\newcommand{\QQ}[0]{\mathbb{Q}}
\newcommand{\RR}[0]{\mathbb{R}}
\newcommand{\Rr}[0]{\mathcal{R}}
\newcommand{\Ss}[0]{\mathcal{S}}
\newcommand{\TT}[0]{\mathbb{T}}
\newcommand{\Uu}[0]{\mathcal{U}}
\newcommand{\Vv}[0]{\mathcal{V}}
\newcommand{\Ww}[0]{\mathcal{W}}
\newcommand{\ZZ}[0]{\mathbb{Z}}
\newcommand{\acts}[0]{\circlearrowright}
\newcommand{\explain}[2] {
	\begin{flalign*}
		 && \text{#2} && \text{#1}
	\end{flalign*}
}
\newcommand{\maketitleprint}[0]{ \begin{center}
	%\begin{tikzpicture}[scale=3]
	%	\duck[graduate=gray!20!black, tassel=red!70!black]
	%\end{tikzpicture}	
	\includegraphics[width=6cm]{cover}
\end{center}
}

% theorem commands
\newtheoremstyle{c_remark}
	{}	% Space above
	{}	% Space below
	{}% Body font
	{}	% Indent amount
	{\bfseries}	% Theorem head font
	{}	% Punctuation after theorem head
	{.5em}	% Space after theorem head
	{\thmname{#1}\thmnumber{ #2}\thmnote{ \normalfont{\text{(#3)}}}}	% head content
\newtheoremstyle{c_definition}
	{3pt}	% Space above
	{3pt}	% Space below
	{}% Body font
	{}	% Indent amount
	{\bfseries}	% Theorem head font
	{}	% Punctuation after theorem head
	{.5em}	% Space after theorem head
	{\thmname{#1}\thmnumber{ #2}\thmnote{ \normalfont{\text{(#3)}}}}	% head content
\newtheoremstyle{c_plain}
	{3pt}	% Space above
	{3pt}	% Space below
	{\itshape}% Body font
	{}	% Indent amount
	{\bfseries}	% Theorem head font
	{}	% Punctuation after theorem head
	{.5em}	% Space after theorem head
	{\thmname{#1}\thmnumber{ #2}\thmnote{ \text{(#3)}}}	% head content

\ifcsname c@english\endcsname
	\theoremstyle{plain}
	\newtheorem{theorem}{Theorem}[section]
	\newtheorem{lemma}[theorem]{Lemma}
	\newtheorem{proposition}[theorem]{Proposition}
	\newtheorem*{proposition*}{Proposition}
	%\newtheorem{corollary}[theorem]{אין חלופה עברית}

	\theoremstyle{definition}
	\newtheorem{definition}[theorem]{Definition}
	\newtheorem*{definition*}{Definition}
	\newtheorem{example}{Example}[section]
	\newtheorem{exercise}{Exercise}[section]

	\theoremstyle{remark}
	\newtheorem*{remark}{Remark}
	\newtheorem*{solution}{Solution}
	\newtheorem{conclusion}[theorem]{Conclusion}
	\newtheorem{notation}[theorem]{Notation}
\else
	\theoremstyle{c_plain}
	\newtheorem{theorem}{משפט}[section]
	\newtheorem{lemma}[theorem]{למה}
	\newtheorem{proposition}[theorem]{טענה}
	\newtheorem*{proposition*}{טענה}
	%\newtheorem{corollary}[theorem]{אין חלופה עברית}

	\theoremstyle{c_definition}
	\newtheorem{definition}[theorem]{הגדרה}
	\newtheorem*{definition*}{הגדרה}
	\newtheorem{example}{דוגמה}[section]
	\newtheorem{exercise}{תרגיל}[section]

	\theoremstyle{c_remark}
	\newtheorem*{remark}{הערה}
	\newtheorem*{solution}{פתרון}
	\newtheorem{conclusion}[theorem]{מסקנה}
	\newtheorem{notation}[theorem]{סימון}
\fi

% Questions related commands
\newcounter{question}
\setcounter{question}{1}
\newcounter{sub_question}
\setcounter{sub_question}{1}

\ifcsname c@english\endcsname
	\newcommand{\question}[1][0]{
		\ifthenelse{#1 = 0}{}{\setcounter{question}{#1}}
		\section{Question \arabic{question}}
		\addtocounter{question}{1}
		\setcounter{sub_question}{1}
	}

	\newcommand{\subquestion}[1][0]{
		\ifthenelse{#1 = 0}{}{\setcounter{sub_question}{#1}}
		\subsection{Part \alph{sub_question}}
		\addtocounter{sub_question}{1}
	}
\else
	\newcommand{\question}[1][0]{
		\ifthenelse{#1 = 0}{}{\setcounter{question}{#1}}
		\section{שאלה \arabic{question}}
		\addtocounter{question}{1}
		\setcounter{sub_question}{1}
	}

	\newcommand{\subquestion}[1][0]{
		\ifthenelse{#1 = 0}{}{\setcounter{sub_question}{#1}}
		\subsection{סעיף \localecounter{letters.gershayim}{sub_question}}
		\addtocounter{sub_question}{1}
	}
\fi

% import lua and start of document
\directlua{common = require ('../common')}

\GetEnv{AUTHOR}

% headers
\author{\AUTHOR}
\date\today

\title{פתרון מטלה 10 --- אנליזה על יריעות, 80426}
% chktex-file 17
% chktex-file 9

\DeclareMathOperator{\vol}{vol}

\begin{document}
\maketitle
\maketitleprint[blue]

\question{}
תהי $M \subseteq \RR^n$ יריעה קומפקטית, ונניח ש־$f : M \to \RR$ חלקה.
יהי $\xi : M \to \RR^n$ שדה וקטורי.

\subquestion{}
נראה שקיימת סביבה פתוחה $U \subseteq \RR^n$ של $M$ והעתקה חלקה $\tilde{f} : U \to \RR$ כך ש־$\tilde{f} |_{M} = f$.
\begin{proof}
	לכל $x \in M$ קיימת פרמטריזציה מקומית של $M$, כלומר $\alpha_x : U_x \to M$ חלקה, ויש לה הרחבה חלקה לקבוצה פתוחה $x \in V_x \subseteq \RR^n$.
	מהפיכות משמאל נסיק שגם $f |_{U_x}$ ניתנת להרחבה, ונסמן $g_x$ כהרחבה זו.
	$M$ קומפקטית ו־$M \subseteq \bigcup_{x \in M} U_x$ ולכן קיים תת־כיסוי סופי $M \subseteq \bigcup_{i = 1}^N U_x$.
	זוהי משפחת קבוצות פתוחות חסומה ולכן יש לה חלוקת יחידה ${\{ h_i \}}_{i = 1}^N$, ונגדיר את הפונקציה $\tilde{f} : U \to \RR$ עבור $U = \bigcup_{i = 1}^N U_{x_i}$ על־ידי,
	\[
		\tilde{f}(x)
		= \sum_{i = 1}^N h_i(x) g_{x_i}(x)
	\]
	ונקבל ש־$\tilde{f} |_M = f$ ישירות מהגדרתה מחלוקת היחידה והשוויון בסביבות המקומיות.
\end{proof}

\subquestion{}
נסיק שקיימת סביבה פתוחה $U \subseteq \RR^n$ של $M$ ושדה וקטורי $\tilde{\xi} : U \to \RR^n$ כך ש־$\tilde{\xi} |_M = \xi$.
\begin{proof}
	נבחין כי קיים פירוק $\xi(x) = (\xi_1(x), \ldots, \xi_n(x))$ כך ש־$\xi_i : M \to \RR$ חלקה לכל $1 \le i \le n$.
	מהסעיף הקודם נובע שקיימת סביבה פתוחה $U_i$ כך ש־$\xi_i$ ניתנת להרחבה חלקה ל־$U_i$.
	עתה נבחר $U = \bigcap_{i = 1}^n U_i$ ונקבל קבוצה פתוחה ב־$\RR^n$ כך ש־$\xi$ ניתנת להרחבה קורדינטה קורדינטה בה לפונקציה $\tilde{\xi} : U \to \RR^n$. \\
	נבחין כי טענה זו נכונה רק עבור מימד סופי.
\end{proof}

\subquestion{}
נראה שקיימות הרחבות $F : \RR^n \to \RR, X : \RR^n \to \RR^n$ המרחיבות את $f, \xi$ בהתאמה.
\begin{proof}
	בסעיף א' הגדרנו חלוקת יחידה של קבוצה פתוחה $M \subseteq U$.
	עתה נגדיר חלוקת יחידה של $\RR^n$ על־ידי הרחבת חלוקת היחידה שהגדרנו בסעיף א', יחד עם הקבוצה $W = \RR^n \setminus C$ עבור $C \subseteq U$ סגורה כך ש־$M \subsetneq C$.
	נגדיר את הפונקציה $F$ על־ידי,
	\[
		F(x)
		= \left( \sum_{i = 1}^N \tilde{f} |_{U_i}(x) \alpha_i(x) \right) + \alpha_W(x) \cdot 0
	\]
	ונקבל פונקציה חלקה כך שהיא מתאפסת מחוץ לסביבה של $U$.
	התהליך עבור $X$ זהה.
\end{proof}

\question{}
\subquestion{}
יהיו $0 < a < b < \pi$,
ונגדיר את היריעה $N_{a, b} = S^{n - 1} \times [a, b] \subseteq \RR^{n + 1}$. \\
נראה ש־$\varphi : N_{a, b} \to S^n$ הנתונה על־ידי $\varphi(x, y) = (x \sin y, \cos y)$ דיפאומורפיזם על תמונתה, כאשר $x \in \RR^{n - 2}, y \in \RR$.
\begin{proof}
	מתקיים,
	\[
		{\lVert \varphi(x, y) \rVert}^2
		= {\lVert (x \sin y, \cos y) \rVert}^2
		= {\lVert x \rVert}^2 \sin^2 y + \cos^2 y
		= 1
	\]
	ישירות מהעובדה כי $x \in S^{n - 2}$.
	לכן $\varphi(x, y) \in S^n$ והפונקציה מוגדרת היטב.
	זוהי גם פונקציה חלקה כהרכבת חלקות, ונותר לבדוק חד־חד ערכיות, ממשפט הפונקציה ההפוכה נסיק דיפאומורפיזם לתמונה. \\
	נניח ש־$x \ne x', y \ne y'$, אם $\cos y \ne \cos y'$ אז סיימנו, אחרת $y = y' = \frac{\pi}{2}$ בלבד (ישירות מתחום ההגדרה).
	במקרה זה נקבל $\sin y = \sin y' = 1$ ולכן $x \sin y \ne x' \sin y'$ וקיבלנו חד־חד ערכיות.
	אם גם $x = -x'$
\end{proof}

\subquestion{}
נסמן $M_{a, b} = \varphi(N_{a, b})$.
נראה שאם $f : S^n \to \RR$ רציפה אז,
\[
	\int_{M_{a, b}} f(x)\ dx
	= \int_{a}^{b} {(\sin y)}^{n - 1} \int_{S^{n - 1}} f(x \sin y, \cos y)\ d\vol_{n - 1}(x)\ dy
\]
\begin{proof}
	ישירות מהגדרת אינטגרל, שימוש ב־$\varphi$ ומשפט פוביני נקבל,
	\[
		\int_{M_{a, b}} f(x)\ dx
		= \int_{a}^{b} \int_{S^{n - 1}} f(x \sin y, \cos y) V(D \varphi |_{(x, y)})\ d\vol_{n - 1}(x)\ dy
	\]
	ולכן מספיק שנראה ש־$V(D \varphi |_{(x, y)}) = {(\sin y)}^{n - 1}$.
	נובע ישירות מחישוב.
\end{proof}

\subquestion{}
נסמן $N = S^{n - 1} \times (0, \pi)$ ונגדיר $\varphi : N \to S^n$ כמקודם.
נראה ש־$S^n \setminus \varphi(N)$ ממידה 0.
\begin{proof}
	נבחין ש־$\varphi(S^{n - 1} \times [0, \pi]) = S^n$ (אילו היינו מרחיבים את ההגדרה) וכן,
	\[
		L
		= \varphi(S^{n - 1} \times \{0, \pi\})
		= \varphi(S^{n - 1} \times \{ 0 \})
		= \{ (x \sin 0, \cos 0) \mid x \in S^{n - 2} \}
		= S^{n - 2} \times \{ 0 \}
	\]
	כלומר $S^n \setminus \varphi(N) = L$ עבור $L$ יריעה ממימד $n - 1$, ובפרט קבוצה ממידה 0.
\end{proof}

\subquestion{}
נסיק שאם $f : S^n \to \RR$ רציפה אז,
\[
	\int_{M_{a, b}} f(x)\ dx
	= \int_{a}^{b} {(\sin y)}^{n - 1} \int_{S^{n - 1}} f(x \sin y, \cos y)\ d\vol_{n - 1}(x)\ dy
\]
\begin{proof}
	מתקבל באופן זהה לחלוטין לסעיף ב'.
\end{proof}

\question[4]
\subquestion{}
תהי $M^k \subseteq \RR^n$ יריעה קומפקטית עם שפה.
נניח ש־$x \in \RR^n \setminus M$ נקודה ו־$p \in M$ כך שמתקיים,
\[
	\lVert p - x \rVert
	= \min_{q \in M} \lVert q - x \rVert
\]
נראה ש־$p - x$ היא אנכית ל־$T_p M$.
\begin{proof}
	אנו רוצים להראות שלכל $v \in T_p M$ מתקיים $\langle v, p - x \rangle = 0$.
	באופן שקול נרצה להראות שלכל מסילה $\gamma : (-\delta, \delta) \to M$ כך ש־$\gamma(0) = p$ מתקיים $\langle \gamma'(0), v - x \rangle = 0$.
	נניח שקיימת מסילה כזו כך שמתקבל ערך שונה מאפס.
	אז מהגדרת מכפלה פנימית קיים $0 < \varepsilon < \delta$ (בלי הגבלת הכלליות) כך ש־$\langle \gamma(\varepsilon), v - x \rangle < \langle \gamma(0), v - x \rangle$, וזו סתירה להגדרת $\gamma(0) = p$.
\end{proof}

\subquestion{}
נמצא דוגמה ליריעה עם שפה $M^k \subseteq \RR^n$ כך שהטענה שהוכחנו זה עתה לא חלה בה.
\begin{solution}
	נבחר את חצי הספירה $\HH^n \cap S^{n + 1}$ ואת הנקודה $x = (1, 0, \ldots, -\epsilon)$ נקודה קרובה מאוד לקוטב.
	בבירור עבור המישור המשיק המורחב נקבל ש־$p - x$ נמצאת במישור המשיק, אבל אם נבחן את המישור המשיק עם התייחסות למשיק חד־צדדי בשפה נקבל שאכן יש אנכיות.
\end{solution}

\subquestion{}


\end{document}
